\documentclass[11pt]{article}
\usepackage{amsmath,amssymb}
\usepackage{geometry}
\usepackage{hyperref}
\geometry{a4paper,margin=1in}
\title{Supplementary Methods: Retrospective Two-Site PV Evaluation}
\author{Research draft — author details require verification}
\date{November 2026}
\begin{document}
\maketitle
\section{Scope}
This supplement documents the empirical analysis in the main manuscript. It does not assert deployment-time validity or finite-sample coverage guarantees under temporal dependence or site shift. Weather inputs are historical estimates.
\section{Interval construction}
For lower and upper quantile estimates $\hat q_{\alpha/2}(x)$ and $\hat q_{1-\alpha/2}(x)$, the signed interval violation score is
\begin{equation}
E_i=\max\{\hat q_{\alpha/2}(x_i)-y_i, y_i-\hat q_{1-\alpha/2}(x_i)\}.
\end{equation}
The implementation computes the higher order statistic on a chronological calibration subset and reports empirical holdout coverage and interval scores. Exchangeability is not established for this temporal/site-shift evaluation, so no guarantee is claimed.
\section{Reproducibility notes}
Environment versions and rerun commands are listed in the project reproducibility report. Dependencies are not fully pinned and LaTeX compilation was unavailable in this environment. Earlier drafts contained unverified parameter searches and environment declarations; those claims are withdrawn.
\end{document}
